% Pursuing Stacks (cotes 134-2 to 134-8): archive cote 134-5, site page 54 (PDF page 55), typescript page 241.
% typescript. Draft compared: ps4.tex:2711-2789.
% AI-transcribed by Claude Opus 5.5 (claude-opus-5-5, adaptive thinking, effort high), Message Batches API, batch ps2 (msgbatch_01NVwr1Yuo8qJf5rykeGBXPK), 2026-11-12; reviewed against both supplied scan halves; see CHECK_DECISIONS.tsv.
\scanpage{134-5/54}
\mspage{241}
taneously for all three.

Having thus well in hand the \xout{all-important}\inserted[typed]{basic} functor (11) (which in
case of example \xout{three} 3) above with $M=A^{\wedge}$, reduces to the all-important
functor $\mathrm{Hot}_{A}\hand{\to}(\mathrm{Hot})$), \xout{Hot} we cannot but define an asphericity structure
to be \hand{\uline{\textcolor{black}{modelizing}}} , as meaning that this canonical functor is an equiva-%
lence of categories. In the case of $M= A^{\wedge}$ above, this means that A
is a pseudo-test category - a relatively weak test notion still. It
appears as just a \xout{bit} small bit stronger than merely assuming the pair
$(M,W)$ to be modelizing, \inserted[typed]{i.e. to be a ``\hand{\uline{\textcolor{black}{modelizer}}}'',} namely assuming $\mathrm{Hot}_{M}$ to be equivalent (in some
way or other ...) to (Hot). Maybe we should be a little more cautious
with the use of the word ``modelizing'' though, and devise a terminology
which should reflect very closely the \unsure{hierarchy|Hierarchy} of \deleted{progressively stronger}
test notions
\[
\text{(pseudo-test cat.)}\hand{\supset}\text{(weak test cat.)}\hand{\supset}\text{(test cat.)}\hand{\supset}\text{(strict test cat.)}\hand{\supset} \inserted{\deleted{\illegible{1}}}\inserted[typed]{(contractors)}
\]
which gradually has peeled out of our earlier reflections, \inserted[typed]{by pinpointing} \xout{with} corres-%
ponding qualifications for an asphericity structure,as \inserted[typed]{being} ``pseudo-modeli\xout{xxx}%
zing'' , \xout{``modelizing''xx} ``weakly modelizing'', ``modelizing'', ``strictly mode%
lizing''\hand{,} and ??? . It now appears that a little extra reflection is nee-%
ded here - for today it's getting a little late though !

18.6\hand{, and 19.6.}

78. It is about time now to get a comprehensive treatment, in the
context of asphericity structures, of the relationships suggested time
ago in the ``observation'' and the \xout{x}``silly question'' of section 46 (pages
94 and 95). The former is concerned with test-functors from test cate-%
gories to modelizers, the second more generally with model-preserving
maps between modelizers, having properties similar to the map $M\hand{\to} A^{\wedge}$
stemming from a test-functor. For the time being, as I found but litt-%
le time for the ``extra reflection'' which seems needed, only the first
situation is by now reasonably clear in my mind.

As no evident \over{p}{s}eries of ``asphericity structure''-notions has appea-%
red, paralleling the series of test-notions \xout{displayed} recalled by the
end of last monday's reflections (see above), I'm going to keep (provi-%
sionally only, maybe) the name of a \hand{\uline{\textcolor{black}{modelizing}}} \hand{\uline{\textcolor{black}{asphericity}}} \hand{\uline{\textcolor{black}{structure}}}
$(M,M_{\mathrm{as}})$ as one for which the canonical functor
\[
(1) \qquad \mathrm{Hot}_{M} = W^{-1}M \hand{\longrightarrow} (\mathrm{Hot})
\]
(cf (14) p. 240) is an equivalence. This notion at any rate is sa-%
\margin{54}
% A small circumflex-like mark is visible above the p of asphericity near the opening paragraph; its function is not securely identifiable.
